\ifdefined\isincluded
\else
\documentclass[a4paper, 12pt]{report}
\usepackage[margin=0.8in]{geometry}
\usepackage{amsmath}
\usepackage{amsfonts}
\usepackage{amssymb}
\usepackage{graphicx}
\usepackage[english]{babel}
\usepackage[autostyle, english = american]{csquotes}
\MakeOuterQuote{"}
\usepackage{setspace}
\onehalfspacing
\usepackage{booktabs}
\usepackage{longtable}
\usepackage{tabularx}
\usepackage{array}
\usepackage{url}
\usepackage{cite}

\makeatletter
\let\oldincludegraphics\includegraphics
\renewcommand{\includegraphics}[2][]{%
  \IfFileExists{#2}{%
    \oldincludegraphics[#1]{#2}%
  }{%
    \IfFileExists{images/#2}{%
      \oldincludegraphics[#1]{images/#2}%
    }{%
      \fbox{\parbox[c][3cm][c]{0.8\linewidth}{\centering \texttt{[Image: #2]}}}%
    }%
  }%
}
\makeatother

\begin{document}
\chapter{Geometric Representation Dynamics and Attractor Basins}
\fi

\label{chap.6}
\section{Motivation and Theoretical Framework}
Understanding the mathematical concepts underpinning latent recursive deliberation is the main scientific objective of this thesis. While traditional agent literature views LLMs as black-box generators, we formulate the internal activations as high-dimensional trajectories on a representation manifold $\mathcal{M} \subset \mathbb{R}^{d_h}$ ($d_h = 2560$).

To evaluate \textbf{Hypothesis H5 (Latent recursions converge into solution attractor basins without dimensional collapse)}, we extract representations $H_{i,r}$ for each of the three agents throughout all 912 iterations and compute seven geometric metrics.

\section{Representation Drift and Convergence Analysis}
Representation drift $D_r$ measures how quickly the latent state changes across subsequent recursive passes:
\begin{equation}
    D_r = \| H_r - H_{r-1} \|_2
\end{equation}
Figure~\ref{fig:drift_decay} illustrates how recursion depth and representation drift are related.

\begin{figure}[h!]
\centering
\includegraphics[width=0.9\linewidth]{drift_decay_vs_recursion.png}
\caption{Latent ReAct's representation drift decay score is 70.Convergence into an attractor basin is confirmed by a 9\% decrease in representation drift ($D_1 = 0.0668 \to D_3 = 0.0194$), even though Text ReAct representation drift flatlines.}
\label{fig:drift_decay}
\end{figure}

As shown in Figure~\ref{fig:drift_decay}:
\begin{itemize}
    \item \textbf{Latent ReAct demonstrates high drift decay ($-70.9\%$):} Drift declines monotonically from $D_1 = 0.0668$ to $D_2 = 0.0211$ and $D_3 = 0.0194$ (decay slope $\approx -0.024$). This illustrates how the recurrent loop compresses state space to direct routes toward a stable fixed-point solution.
    \item \textbf{Text ReAct flatlines ($-9.0\%$):} There is relatively little drift decay in the reconstructed internal activations in Text ReAct ($D_1 = 0.0402 \to D_3 = 0.0366$, slope $\approx -0.002$). Without latent recurrent feedback, text-mediated reasoning explores the manifold without stabilizing.
\end{itemize}

\section{Manifold Effective Rank and Dimensional Stability}
To determine if recursive latent feedback causes representation collapse, in which all vectors collapse into a straightforward low-dimensional subspace, we compute the \textbf{Effective Rank ($\text{ER}$)} \cite{roy2007effective} of the representation covariance matrix:
\begin{equation}
    \text{ER}(H) = \exp \left( -\sum_{k=1}^K p_k \ln p_k \right), \quad p_k = \frac{\sigma_k}{\sum_{j} \sigma_j}
\end{equation}
where $\sigma_k$ represents the singular values of the centered representation matrix $H$.

\begin{figure}[h!]
\centering
\includegraphics[width=0.85\linewidth]{effective_rank_spectrum.png}
\caption{Effective Rank Spectrum: The stability of the Effective Rank over recursion rounds ($r=0 \to r=3$) confirms zero dimensional collapse.}
\end{figure}

As shown in Figure~\ref{fig:effective_rank}:
\begin{itemize}
    \item Latent ReAct maintains an effective rank of $\text{ER} \approx 94.45$ at $r=1$ during active error correction, peaks at $r=2$ ($\text{ER} = 94.70$), and settles at $\text{ER} = 92.99$ at $r=3$.
    \item By retaining rich, non-degenerate geometric expressiveness throughout all recursive rounds, the manifold conclusively refutes the idea of dimensional collapse.
\end{itemize}

\section{Inter-Agent Representation Distance and Specialization}
To ascertain whether the sub-agents maintain functional role specialization or merge into homogenous representations, we compute the inter-agent cosine similarity $\text{CosSim}(A_i, A_j)$ and Euclidean distance $D_{ij} = \|H_i - H_j\|_2$.

\begin{figure}[h!]
\centering
\includegraphics[width=0.85\linewidth]{inter_agent_cosine_heatmap.png}
\caption{Inter-Agent Representation Cosine Similarity: Near-orthogonal agent specialization in Latent ReAct ($\text{CosSim} < 0.087$) against redundant prompt overlap in Text ReAct ($\text{CosSim} > 0.998$).} \label{fig:cosine_heatmap}
\end{figure}

As depicted in Figure~\ref{fig:cosine_heatmap}:
\begin{itemize}
    \item With $\text{CosSim}(A_1, A_2) = 0.087$ and $\text{CosSim}(A_2, A_3) = 0.049$, the agents in \textbf{Latent ReAct} are almost orthogonal. The average distance between agents is $D_{ij} \approx 51.5$. Each agent inhabits a unique subspace suitable for its role.
    \item Internal states in \textbf{Text ReAct} substantially overlap ($\text{CosSim} > 0.998$, $D_{ij} \approx 10.6$). Verbal instructions are insufficient to achieve internal representational specialization.
\end{itemize}

\section{PCA Phase-Space Trajectory Flow and Attractor Basins}
To observe trajectory progression in phase space, we use Principal Component Analysis (PCA) to all 912 execution trajectories, projecting 2560-dimensional states onto the first two principal components ($PC_1, PC_2$).

\begin{figure}[h!]
\centering
\includegraphics[width=0.9\linewidth]{pca_trajectory_attractors.png}
\caption{PCA Phase-Space Trajectory Flow: Streamlines demonstrate the dynamic contraction of divergent initial states into a compact solution attractor basin centered at $(PC_1 \approx -12.5, PC_2 \approx +2.5)$.}
\end{figure}

Figure~\ref{fig:pca_attractors} reveals clear phase-space dynamics:
\begin{enumerate}
    \item \textbf{Exploration Phase ($r = 0 \to 1$):} The significant variety of the trajectories over the peripheral ($PC_1 \in [-30, +10]$) reflects different initial issue conditions.
    \item \textbf{Contraction Stage ($r = 1 \to 2$):} When streamline vector fields converge inward, divergent trajectories are realigned.
    \item \textbf{Attractor Basin ($r = 3$):} Successful trajectories converge firmly in an invariant attractor basin centered at $(PC_1 \approx -12.5, PC_2 \approx +2.5)$. Conversely, failure routes diverge into outer orbits.
\end{enumerate}

\section{Subspace Angles and Centroid Separability}
Using SVD decomposition, we determine the major subspace angle $\theta_1$ and centroid distance $\Delta_{\text{SF}}$ between successful and unsuccessful trajectories:
\begin{itemize}
    \item \textbf{Principal Subspace Angle:} $\theta_1 = 19.83^\circ$ (up to $70.66^\circ$ for complex projects), demonstrating a notable angular divergence between the success and failure manifolds.
    \item \textbf{Centroid Separability:} $\Delta_{\text{SF}} = 0.7617$ shows that effective solutions occupy a unique, separable geometric cluster on the manifold.
\end{itemize}
These empirical metrics support Hypothesis H5 by showing that latent recursion functions as a geometric dynamical system driven by attractor convergence.

\ifdefined\isincluded
\else
\begin{thebibliography}{99}
\bibitem{roy2007effective} O.~Roy and M.~Vetterli, ``The effective rank: A measure of effective dimensionality,'' in \emph{15th European Signal Processing Conference (EUSIPCO)}, pp.~606--610, 2007.
\end{thebibliography}
\end{document}
\fi
